% ED26-SRC-063-0105: Ganeri, Indian Logic, HHL 1 (2004), pp.309--395.
% Source version: held publisher-typeset chapter, not the distinct 1913 Nyaya text.
\chapter{Indian Logic: Debate, Reasons, Standpoints, and Absence}
\label{sec:p13-ganeri-indian-logic}
\chaptermark{Indian Logic}
\index{Indian logic}\index{Ganeri, Jonardon}
\index{Nyaya}\index{Buddhist logic}\index{Jaina logic}\index{Navya-Nyaya}

Indian logical traditions do not share a single object language or a single
criterion of a good argument. Jonardon Ganeri's 2004 chapter follows four
connected problems: the regulation of debate; the connection between a reason
and what it establishes; assertion from different standpoints; and the logical
analysis of properties, absence and number
\cite[pp.~309--391]{Ganeri2004IndianLogicPublished}. A disputant's permissible move,
an epistemically reliable inference, a standpoint-relative assertion and an
ontological dependence are different objects. A formal treatment of one must
not be credited automatically to the others.

The historical formulations below are reported through Ganeri's chapter.
His interpretations using natural deduction, sets, modal or assertion logic,
and directed graphs are explicitly modern reconstructions. They are not
ancient Sanskrit rule sheets. The local mathematical results establish the
displayed reconstructions under their stated assumptions, not the adequacy
of every historical interpretation. This presentation is distinct from the
1913 English translation of the \emph{Nyaya-sutra}.

\section{An argument is also an exchange}
\label{sec:p13-ganeri-debate}

\subsection{Question, commitment, and the right to retract}

Ganeri begins with interrogation rather than a truth-functional calculus:
the exchanges involving Janaka and Yajnavalkya in the Upanishads, and the
dialogues of Milinda and Nagasena. Questions can test an interlocutor's
understanding and expose a commitment, not merely request information.
The contrast between scholarly discussion and royal coercion matters:
a scholar can acknowledge an error and revise a view; punishment of dissent
changes the exchange's aim and the significance of assent
\cite[pp.~309--314]{Ganeri2004IndianLogicPublished}.

A yes-or-no demand may conceal a presupposition. Refusal to answer need not
mean ignorance; it can reject the question's framing. Answers admitting both
alternatives, neither alternative, or a qualified distinction do not by
themselves introduce a four-valued semantics. The lamp whose successive flames
are causally connected, and the child who becomes an adult, distinguish
continuity from strict numerical identity and from complete unrelatedness.
Likewise, memory need not be classified exclusively as an internal operation
or an external aid. These cases concern the adequacy of the offered alternatives.

\subsection{The Kathavatthu: a reason to retract is not an axiom}

The \emph{Kathavatthu} organizes disputes by eight openings: an unqualified
affirmative and negative, followed by their qualifications concerning place,
time, and the range of things. The five phases distinguish a forward exchange,
a reverse exchange, refutation, application, and conclusion. A counter-question
can return a burden of explanation; a defensive argument can attack the
initial reason. These are not interchangeable occurrences of implication
\cite[pp.~314--321]{Ganeri2004IndianLogicPublished}.

In the dispute about a person or soul as an ultimate reality, let \(P\) mean
that it is known as a real ultimate, and \(Q\) that it is known in the same way
as other ultimate realities. A respondent accepts \(P\), rejects \(Q\), and is
challenged about their connection. In the reconstruction attributed to
Matilal, the additional premise \(P\to Q\), with \(\neg Q\), permits rejection
of \(P\). Bochenski also uses contraposition and modus tollens.
Ganeri's alternative treats the consequence \(Q\) as extracted from the
respondent's commitments: the conditional can summarize that extraction
instead of being a separately conceded material implication. The question
is then which commitment the respondent should retract.

\begin{hblemma}{thm:p13-ganeri-bridge}
 {proved=annex:p13.ganeri.bridge}{What an explicit inferential bridge supplies}
In classical first-order reconstruction,
\(\forall x(Fx\to Gx),Fa\) entails \(Ga\), and \(P\to Q,\neg Q\)
entails \(\neg P\). The three facts \(Fa,Fb,Gb\) alone do not entail \(Ga\).
Consequently, an accepted example is not by itself a universal bridge.
\end{hblemma}

External negation of a whole assertion must also be distinguished from
predication of a negative property. Rejecting that \(c\) is \(D\) does not
automatically endorse that \(c\) instantiates an independently defined
negative kind. The classical derivation just given is an interpretation of
an argument, not proof that every dialogue participant accepts classical
negation or the same universal premise.

\subsection{Nyaya: five members and the relevance of an example}

The \emph{Nyaya-sutra} account joins knowledge-producing means, supporting
reasoning, and regulated disputation. Honest debate, \emph{vada}, is oriented
toward establishing and refuting by accepted means of knowledge and supporting
argument. Competitive \emph{jalpa} admits obstructive devices; \emph{vitanda}
attacks a thesis without establishing a counterthesis. A refutation-only
strategy is not therefore identical with a demonstration of the opposite
\cite[pp.~321--334]{Ganeri2004IndianLogicPublished}.

The five members are \emph{pratijna}, the thesis; \emph{hetu}, the reason;
\emph{udaharana}, the example; \emph{upanaya}, its application to the subject;
and \emph{nigamana}, the conclusion. In a schematic sound argument these become:
sound is transitory; it is produced; a pot supplies a produced transitory
case; sound has the relevant feature; hence the thesis is reaffirmed.
The example's role must be specified. An early case-to-case reading emphasizes
relevant similarity, while later readings make the reason--target relation
explicit. Schayer's modern natural-deduction reading supplies
\(\forall x(Fx\to Gx)\) and \(Fa\), instantiates at \(a\), and concludes
\(Ga\). It is not a derivation of \(Fa\) from \(Ga\).

Similarity without relevance is a source of faulty comparison. Four legs
do not isolate cows from horses; a dewlap is offered as a more discriminating
feature within the assumed comparison. An inference that the soul is eternal
because it is uncreated like space faces an opposed comparison based on
perceptibility and pots. Accepting both analogies without examining their
connecting relations leaves the disagreement unresolved.

Ganeri's case-based illustration concerns two mushroom eaters and a matching
cluster of symptoms. The comparison can guide a reasoner who does not know
the common underlying mechanism; its support depends on the relevance and
stability of the matched features. It is a schematic epistemological example,
not medical advice. A cooked grain is a useful sample of the pot's rice only
given adequate common cooking conditions; a drop is a useful sample of sea
water only with appropriate homogeneity. Sampling does not establish those
conditions by its own name. Ants carrying eggs, a full river, and smoke
likewise illustrate different explanatory connections to future rain,
past rain, and fire, not one unrestricted resemblance principle
\cite[pp.~324--329]{Ganeri2004IndianLogicPublished}.

\subsection{Four alternatives, presuppositions, and tarka}

Nagarjuna's fourfold challenges and the later contrast between independent
inference and consequence-based refutation concern what a position
presupposes. Ganeri illustrates a relational fourfold division. For a set
\(D\), \(a\in D\), and \(R(a)=\{b\in D:R(a,b)\}\), the possibilities are:
\(R(a)=\{a\}\); nonempty \(R(a)\) excluding \(a\); \(R(a)\) containing \(a\)
and another element; or empty \(R(a)\).

\begin{hblemma}{thm:p13-ganeri-fourfold}
 {proved=annex:p13.ganeri.fourfold}{A relational fourfold division and its limit}
The four displayed possibilities are disjoint and exhaustive. Over the real
numbers, with \(R(a,b)\) meaning \(b^2=a\), they are realized respectively
at \(a=0,4,1,-1\). An equivalence relation need not exclude the third case.
\end{hblemma}

This partition does not establish that a particular historical fourfold
argument has that semantics. Moreover, the last sentence corrects the
restriction asserted in Ganeri's graph-independent example:
an equivalence class may contain its member and another member.

\emph{Tarka} supports inquiry by considering consequences, challenging an
assumption, or shifting a burden; it is not simply another independent
knowledge-producing means. Seeing something resembling a person or a stump
leaves a perceptual ambiguity; contextual considerations, such as an aligned
series of objects, can alter a presumption without changing the original
perceptual content. Udayana's inadmissible supports include self-dependence,
mutual dependence, a circle, an ungrounded regress, and an unacceptable
consequence. Sriharsa broadens the challenge to include self-contradictory
speech acts, epistemic self-defeat, equally faulty parallel reasoning, and
lack of discriminating evidence. Neither a counter-argument nor a simplicity
preference alone proves the contrary thesis
\cite[pp.~334--340]{Ganeri2004IndianLogicPublished}.

An \emph{upadhi} can disclose an omitted condition. In the familiar
smoke--fire discussion, reversing a smoke-to-fire connection requires an
additional condition concerning the fuel. This is an assumed explanatory
case, not a universal empirical law about every modern instance of combustion.
Ganeri also gives an abstract characterization. Write \(F\preceq G\) for a
reflexive transitive relation of pervasion:
\[
 F\preceq G\quad\Longleftrightarrow\quad
 \forall Z\,(G\preceq Z\ \Longrightarrow\ F\preceq Z).
\]
\begin{hblemma}{thm:p13-ganeri-pervasion}
 {proved=annex:p13.ganeri.pervasion}{The abstract pervasion characterization}
The displayed equivalence holds in every preorder. It specifies no empirical
method for establishing that a particular smoke--fire relation belongs
to that preorder.
\end{hblemma}

\section{Buddhist and Jaina tests of a reason}
\label{sec:p13-ganeri-reasons}

\subsection{Three marks and what counts as a like case}

Di\.nn\={a}ga's account requires a reason to attach to the subject, to occur in
like cases, and not to occur in unlike cases. The comparison concerns the
property to be established, not merely the reason itself
\cite[pp.~340--344]{Ganeri2004IndianLogicPublished}.
Let \(p\) be the site, \(H\) the reason and \(S\) the target.
The first condition requires \(H(p)\). The other conditions depend on how
the positive comparison class \(\mathsf{sp}\) and negative class
\(\mathsf{vp}\) are fixed. A relational reading can compare objects' states
without first deciding that the target holds at the disputed site.
A class reading excluding the site from positive examples avoids directly
assuming the thesis, but leaves a gap between tested comparison cases
and the site. Neither reading should be replaced silently by the other.

Restricting the reason to like cases is not the same as requiring it in
every like case. \(H\subseteq S\) excludes reason-bearing counterexamples;
it does not say \(S\subseteq H\). A positive instance can be a separate
non-vacuity condition. The formal distribution of properties and the means
of knowing that distribution are separate questions.

Consider the produced-sound case with the disputed sound \(p\), a pot \(a\),
and space \(b\). The commitments used by the example are \(H(p)\),
\(H(a)\land S(a)\), and \(\neg H(b)\land\neg S(b)\), where \(H\) is
produced and \(S\) transitory. These place the comparisons in cell 2;
they do not settle the reason's global distribution by examining only
these two objects. If a separate necessity premise
\(\forall z(Hz\to Sz)\) is known, the bridge lemma gives \(S(p)\).
By contrast, when \(H\) is transitory and \(S\) man-made, lightning supplies
\(H(c)\land\neg S(c)\) within the example's assumed ontology.
This is a concrete failure of the required bridge, not merely a less
persuasive similarity.

\subsection{Nine reasons, empty classes, and counterexamples}

For a nonempty comparison class \(C\), use \(+\) when every member has \(H\),
\(-\) when none does, and \(\pm\) when some but not all do.
Ganeri's ninefold wheel has positive-class status as its row and
negative-class status as its column. In row-major order the cells are
\(1,2,3;4,5,6;7,8,9\). Its good reasons are 2 and 8:
positive examples exist and the reason is absent from the negative class
\cite[pp.~344--350, figures~2--4]{Ganeri2004IndianLogicPublished}.

The wheel's sound examples make the distinctions concrete. A reason based
on knowability for eternity is too broad (1). Produced sound inferred
transitory has a pot as a positive comparison and space as a negative
comparison (2). Inferring man-made from transitory fails at lightning (3).
Produced offered for eternal reverses the relevant connection (4).
Audibility offered for eternity is uniquely confined to the disputed sound
class and supplies no independent positive comparison (5). Man-made offered
for eternal has a produced pot as an adverse instance (6).
Transitory offered for natural encounters the contrast between lightning
and a pot (7). Man-made offered for transitory has positive pots without
a reason-bearing negative instance in the assumed range (8).
Incorporeality offered for eternal crosses the contrast among space,
pots and action (9). These are the chapter's doctrinal classifications,
not newly verified empirical assertions about all their subjects.

\begin{hblemma}{thm:p13-ganeri-wheel}
 {proved=annex:p13.ganeri.wheel}{The finite wheel as a distribution test}
For nonempty disjoint classes \(\mathsf{sp},\mathsf{vp}\), the three statuses
in each coordinate give exactly nine possibilities. Positive occurrence
and absence in \(\mathsf{vp}\) select precisely cells 2 and 8.
If the site is in neither class, passing this test alone does not imply
\(S(p)\), even when \(H(p)\).
\end{hblemma}

An empty class needs a fourth status \(0\): otherwise universal truth and
absence of any witness are conflated. Uddyotakara's expanded wheel uses
\(+,-,\pm,0\) on both axes. In Ganeri's figure the accepted cells are
\((+,-),(+,0),(\pm,-),(\pm,0),(0,-)\). They are not obtained merely
by adding empty classes to the positive-witness condition of the ninefold
wheel. A universally applicable reason may have no negative comparison,
while a distinguishing negative pattern may have no separate positive
comparison. What serves as a site or a comparison range must therefore
be stated, rather than imposed from a modern point-subject model.

\subsection{Why a connection holds}

Dharmakirti distinguishes an effect, an essential property, and appropriate
non-observation as grounds of inference. Smoke as an effect of fire is not
just a frequently matching appearance. A \emph{simsapa} as a tree appeals
to an essential connection. Ganeri's token-identity interpretation of
property instances is a modern analysis of that connection, not a historical
set-theoretic definition. Taste and colour can be associated through a
common cause; day and night can be explained through an intervening
mechanism. Transferring an effect through an intermediate cause requires
the needed necessary connection at each step
\cite[pp.~350--352]{Ganeri2004IndianLogicPublished}.

Non-observation supports absence only with a perceptibility premise:
if the object were present under suitable observational conditions,
it would be encountered. Merely failing to look supplies no such premise.
The logical condition for a good reason does not itself establish
causal necessity, conceptual inclusion, or adequate observation.

The inferential form makes the observational condition explicit.
Let \(C\) assert that the relevant conditions for observation obtain,
\(E\) that the object is present, and \(O\) that it is observed.
From \(C\), \(C\land E\to O\), and \(\neg O\), an assumption of \(E\)
gives \(C\land E\), then \(O\), a contradiction, hence \(\neg E\).
If \(C\) is omitted, the assignment \(C=0,E=1,O=0\) satisfies
the conditional and non-observation without absence. This is the
classical reconstruction of the limiting premise, not an unrestricted
historical rule that every unobserved object is nonexistent.

\subsection{Jaina necessitation and different sites}

Jaina authors challenge the identification of a property's absence with
its mere non-presence. Under an account admitting presence and absence
together, excluding the reason wherever the target is absent is not the
same as excluding reason without target. The single requirement expressed
by \emph{avinabhava} or \emph{anyathanupapannatva} is that the reason cannot
occur otherwise than with what it establishes
\cite[pp.~352--354]{Ganeri2004IndianLogicPublished}.

Examples involving a star's rise or sunlight invite a common temporal site,
but the moon and its reflection involve distinct objects. Ganeri therefore
considers a cross-site reconstruction:
\[
 \forall x\exists y(Fx\to Gy),\qquad \exists xFx
 \quad\Longrightarrow\quad \exists yGy.
\]
\begin{hblemma}{thm:p13-ganeri-crosssite}
 {proved=annex:p13.ganeri.crosssite}{A reason and conclusion at different sites}
The displayed inference is valid in classical first-order semantics.
It does not imply that the object satisfying \(F\) itself satisfies \(G\).
\end{hblemma}
The existentially bound second site is essential. The modern formula
does not claim that all historical Jaina reasons were expressed in a
quantified calculus.

\section{Jaina standpoints: assertion is not absolute truth}
\label{sec:p13-ganeri-standpoints}

\subsection{Seven qualified predications}

A \emph{naya} is a perspective within which assertions can be consistent.
Claims belonging to different perspectives need not form one consistent
unqualified theory. The seven \emph{syat} predications distinguish
affirmation, denial, their successive combination, non-assertibility,
affirmation with non-assertibility, denial with non-assertibility, and the
combination of all three. The historical account's simultaneous
non-assertibility is not automatically an asserted contradiction
\cite[pp.~354--363]{Ganeri2004IndianLogicPublished}.

Ganeri compares rival interpretations of the third component.
The conjunction interpretation associated with Bharucha and Kamat, and
the joint-assertion reading discussed by Matilal, are distinct proposals.
If non-assertibility means assertion of \(p\land\neg p\), or both assertions
within one standpoint, suitable conjunction principles collapse distinctions
that the sevenfold scheme was meant to preserve. His alternative allows
a standpoint to assert neither \(p\) nor \(\neg p\).
Write \(\sigma:p\) for assertion at a standpoint and
\(\nabla p\) for assertion at some standpoint in a nonempty discourse
\(\Sigma\). The three support markers are
\[
 t:\exists\sigma\,(\sigma:p),\qquad
 f:\exists\sigma\,(\sigma:\neg p),\qquad
 u:\exists\sigma\,(\neg(\sigma:p)\land\neg(\sigma:\neg p)).
\]
They are not three exclusive truth values of the underlying proposition.
They record which local assertion states occur across a discourse.
Under consistent standpoints, each local state is affirmative, negative,
or neutral. The seven nonempty subsets of these three states give the
seven discourse patterns.

\begin{hblemma}{thm:p13-ganeri-standpoints}
 {proved=annex:p13.ganeri.standpoints}{Seven support patterns without global adjunction}
In the standpoint model of consistent classical theories, every nonempty
support subset of \(\{t,f,u\}\) is realizable. Assertion of \(p\) somewhere
and assertion of \(q\) somewhere do not entail assertion of \(p\land q\)
somewhere. At a single standpoint, conjunction closure does hold.
\end{hblemma}

For example, one standpoint can contain \(p\) and another \(\neg p\)
without either standpoint containing a contradiction. This is a difference
between an inconsistent union and an internally inconsistent standpoint.
Neutrality is local failure to assert either alternative, not a finding
that the proposition lacks any truth. Nor does the existence of several
arguable positions prove that they exhaust reality.

\subsection{Correlation and the limit of a truth table}

When each standpoint is complete, an atom has one of three discourse
patterns: affirmed everywhere, denied everywhere, or affirmed in some
and denied in others. Number these patterns 1, 2, and 3 respectively.
If both \(p\) and \(q\) have the mixed pattern,
the pattern of their conjunction depends on whether their affirmative
standpoints coincide. It can be denied everywhere or mixed.

\begin{hblemma}{thm:p13-ganeri-correlation}
 {proved=annex:p13.ganeri.correlation}{Equal marginal patterns, different conjunctions}
Two finite families of complete classical standpoints can assign the
same mixed pattern to both \(p\) and \(q\), while assigning different
patterns to \(p\land q\). Therefore these marginal patterns alone do not
determine a single-valued conjunction table.
\end{hblemma}

In the three-pattern numbering just used, the two possible outputs for
mixed inputs are 2 and 3: denied everywhere or mixed. They are not
1 and 2; the worked valuation families in the annex fix the range directly.

Ganeri discusses quasi-truth-functions and Jaskowski-style modal comparisons,
but does not complete a seven-valued compositional semantics. A proposed
connection between truth and somebody's assertion is an additional
assumption, not a consequence of standpoint plurality.
His Rescher-inspired assertion kernel has nonvacuity, same-standpoint
conjunction, no asserted contradiction, and closure under consequence.
The stronger condition excluding neutral standpoints is optional.
Rule sheet \ref{sheet:p13-ganeri-assertion-kernel} records this as a modern kernel; it is not a
complete historical Jaina calculus or a proved modal axiomatization
\cite[pp.~363--369]{Ganeri2004IndianLogicPublished}.

\section{Navya-Nyaya: dependencies, absences, and collections}
\label{sec:p13-ganeri-navya}

\subsection{A graph of inherence is an interpretation}

The classical categories include substance, quality, motion, universal,
individuator, and inherence; absence is subsequently added. Ganeri reads
inherence as directed dependence, with an edge pointing from what inheres
to its bearer. Qualities and motions can have a zero-indegree universal
as a predecessor and a substance as a successor. Universals and individuators
are distinguished by their outgoing valencies. Substance compounds can
depend on substances. This is a graph-theoretic proposal for articulating
ontological structure, not a claim that historical authors drew these graphs
\cite[pp.~370--377]{Ganeri2004IndianLogicPublished}.

The minimal graph vocabulary does not distinguish quality from motion:
additional contact or causal relations, or an ontological identification,
are needed. Nor does directedness alone rule out circles or guarantee
a longest incoming path. The following explicit finite model supplies
the assumptions required by the level argument. Its node types are
\(\mathsf U,\mathsf I,\mathsf Q,\mathsf M,\mathsf S\).
Universals and individuators have indegree zero; every quality or motion
has a predecessor in \(\mathsf U\cup\mathsf I\), points only to substances,
and has no other kind of predecessor. Substances have predecessors from
\(\mathsf Q\cup\mathsf M\cup\mathsf S\); every substance is reached from
a quality or motion through zero or more substance nodes.
All substance-to-substance edges belong to a strict transitive
proper-part relation. No other edges occur.

\begin{hblemma}{thm:p13-ganeri-graph}
 {proved=annex:p13.ganeri.graph}{A well-founded finite typed dependence model}
Under these explicit assumptions the graph is acyclic. The longest incoming
path defines a finite level: universals and individuators have level zero,
qualities and motions level one, and substances level at least two.
Removing the strict-part condition permits a substance cycle; removing
finiteness permits unbounded incoming path lengths.
\end{hblemma}

The qualifications matter. A three-node substance cycle is compatible with
the bare statement that only substances inhere in substances. The
strict-part condition excludes it. Finite acyclic dependence gives the
needed maximum; infinite acyclic dependence need not do so. These repaired
assumptions make the local construction usable without attributing its
unqualified conclusion to the historical graph definitions.

\subsection{Absence is not an untyped complement edge}

An absence has an absentee, a location, and relevant temporal and
delimiting conditions. It is not merely a missing edge in the positive
inherence graph. Ganeri's extension distinguishes the link to the absentee
from the link locating the absence. Difference is treated differently
again: an object's otherness concerns other objects, not just an absence
at one chosen location. Reusing the same untyped arrow for all these
relations would erase the proposal's distinctions
\cite[pp.~377--385]{Ganeri2004IndianLogicPublished}.

An absence of the absence must retain its own relational type.
Raghunatha's treatment, as interpreted here, need not identify an object's
double absence with the positive object itself. Ganeri instead treats
the double absence as connected by absence-specific links.
His identification of third and first absence depends on uniqueness
for absences with matching absentee and locating relations.
Agreement of truth indicators alone does not prove identity of these
ontological nodes. The truth-pair model below proves its own periodicity,
not this stronger ontological uniqueness claim.

Presence and absence can coexist under different delimiting conditions:
a monkey contacts one portion of a tree without contacting another;
a pot may be painted on its outside and not on its inside.
An account must specify what the predication and its absence range over.
Classical explosion cannot be imported merely because the word
negation has been used.

In Ganeri's higher-absence reconstruction, external denial of truth and
object-level absence remain distinct. Here \(T\alpha\) states that the
object-language expression \(\alpha\) is true at the specified location;
\(\neg T\alpha\) externally denies that statement. In contrast,
\(\neg\alpha\) forms an object-level absence expression and is not
defined as the Boolean complement of \(T\alpha\).
His displayed constraints are
\[
 \neg T\alpha\ \Longrightarrow\ T\neg\alpha,\qquad
 T\neg\neg\alpha\ \Longleftrightarrow\ \neg T\neg\alpha.
\]
The converse \(T\neg\alpha\Rightarrow\neg T\alpha\) is not accepted.
The typed constraints are collected in sheet \ref{sheet:p13-ganeri-higher-absence}.
Thus double-negation elimination is compatible with failure of
double-negation introduction.

% Ganeri2004 modern reconstruction, pp.367--368: assertion kernel.
% NOT a complete historical Jaina calculus or a complete modal axiomatization.
\calculussheet{sheet:p13-ganeri-assertion-kernel}
 {Ganeri's modern standpoint-assertion kernel}
 {Formula metavariables r,s; standpoint sigma; nonempty discourse Sigma.
 Assertion is typed separately from object-level truth and classical consequence.
 The optional completeness condition is not part of the basic kernel.}
\hbrule{rule:p13-ganeri-nonvacuity}{A1: nonvacuity}
 {\sigma\in\Sigma}{(\exists r)\,\sigma:r}
 {The quantifier ranges over formulas in the metalanguage, not over individuals.}
\hbrule{rule:p13-ganeri-same-conjunction}{A2: same-standpoint conjunction}
 {\sigma:r\quad\sigma:s}{\sigma:(r\land s)}
 {Both premises concern the same standpoint \(\sigma\).}
\hbrule{rule:p13-ganeri-consistency}{A3: standpoint consistency}
 {\sigma\in\Sigma}{\neg(\sigma:(r\land\neg r))}
 {The outer negation denies assertion; the inner negation is object-language negation.}
\hbrule{rule:p13-ganeri-consequence}{R: assertion closure}
 {r\vdash_{\mathrm{CL}}s\quad\sigma:r}{\sigma:s}
 {The consequence relation is the chosen classical reconstruction; this is not a rule
 for combining assertions from different standpoints.}
\hbrule{rule:p13-ganeri-some-assertion}{Some-standpoint assertion}
 {\sigma\in\Sigma\quad\sigma:r}{\nabla r}
 {By definition \(\nabla r\) abbreviates \(\exists\sigma\in\Sigma\,(\sigma:r)\).}
\hbrule{rule:p13-ganeri-assertion-witness}{Unpack some-standpoint assertion}
 {\nabla r}{\exists\sigma\in\Sigma\,(\sigma:r)}
 {This is abbreviation expansion, not an additional historical axiom.}
\hbrule{rule:p13-ganeri-no-neutral}{A4: optional exclusion of neutral standpoints}
 {\sigma\in\Sigma}{(\sigma:r)\lor(\sigma:\neg r)}
 {Optional stronger reconstruction only; it is not valid for arbitrary consistent
 incomplete standpoints.}

% Ganeri2004 p.385: constraints for his modern higher-absence analysis.
\calculussheet{sheet:p13-ganeri-higher-absence}
 {Ganeri's two higher-absence truth constraints}
 {Object-level absence negation is distinct from external denial of T.
 These constraints do not supply a complete binary-connective semantics.}
\hbrule{rule:p13-ganeri-absence-r1}{R1}
 {\neg T\alpha}{T\neg\alpha}
 {The premise is external denial of truth, not an asserted object-level absence.}
\hbrule{rule:p13-ganeri-absence-r2}{R2}
 {}{T\neg\neg\alpha\leftrightarrow\neg T\neg\alpha}
 {The rightmost negation is external; the two negations before \(\alpha\) are object-level.}


\begin{hblemma}{thm:p13-ganeri-absence}
 {proved=annex:p13.ganeri.absence}{A local truth-pair model for higher absence}
Let the initial pair of truth indicators for \(p,\neg p\) be
\((0,1),(1,0)\), or \((1,1)\), and for \(n\geq1\) set
\(T(\neg^{n+1}p)=1-T(\neg^n p)\).
This model satisfies the two displayed constraints at every negation depth.
It validates \(\neg\neg p\Rightarrow p\), but need not validate
\(p\Rightarrow\neg\neg p\).
\end{hblemma}

This model concerns the stated constraints, not a complete semantics for
all Navya-Nyaya connectives. Truth preservation from \(\alpha\) to \(\beta\)
at every location is distinct from requiring the absence of \(\beta\)
to preserve the absence of \(\alpha\). Co-location can separate them.
Gangesa's successive proposed definitions of pervasion refine which
locations, absences and delimiting relations are relevant; a maximal
property with no appropriate absence exposes a definition's applicability
limit. Ganeri's presentation of the final strategy uses a contraposed
property-identity condition rather than treating an undefined absence
as an ordinary empty extension
\cite[pp.~385--387]{Ganeri2004IndianLogicPublished}.

The five provisional tests concern: exclusion of the reason from a
target-absence location; that exclusion restricted to locations different
from target-bearing ones; non-coincidence with difference from a target
location; an absence of the reason found at every target-absence location;
and exclusion from whatever is other than a target location.
These formulations make different uses of absence and otherness.
A target present everywhere defeats their common assumption that
its absence is available with an exemplifying location.
The preferred formulation asks whether any property whose absence
coincides with the reason can be identical to the target. It thereby
avoids directly designating the problematic target-absence.
When presence and absence can coexist, the second provisional test
also avoids forbidding a legitimate reason at a location where the
target is both present and absent. It forbids reason without target,
not simply reason with target-absence.

\subsection{Number: unique possession and collective possession}

The treatment of number raises another typing distinction. Classical
accounts locate numerical qualities and universals in their bearers;
Raghunatha's separate category of number permits counting across
categories. A collection of two pots does not make each pot individually
two. Collective \emph{paryapti} must not be replaced by distributive
inherence \cite[pp.~387--391]{Ganeri2004IndianLogicPublished}.

For Gadadhara's analysis of unique possession of \(F\), Ganeri supplies the
equivalent forms
\[
 Fx\land\neg\exists y(Fy\land y\ne x),\qquad
 Fx\land\forall y(Fy\to y=x).
\]
\begin{hblemma}{thm:p13-ganeri-number}
 {proved=annex:p13.ganeri.number}{Unique possession and a finite collective model}
The two formulas are equivalent in classical first-order logic with equality.
In a finite collective model, an edge from a numerical node \(n\) to a set
\(C\) can mean \(|C|=n\); it does not entail that every member of \(C\)
has the numerical property \(n\).
\end{hblemma}

The collective target is a set of distinct nodes, not a tuple that counts
one node twice. Number nodes may themselves be counted if the model allows
them as targets, but this observation establishes no unrestricted existence
principle for arbitrary collections and no general non-well-founded
set theory. Ganeri ends with a comparative programme: formal interpretations
can reveal inferential structure, provided their added assumptions and
their historical scope remain visible.
